% P8 compliance grounding: SAR narratives (FinCEN) + adverse action (ECOA/Reg B)
\section{Compliance Grounding for the Scribe}
\label{sec:p8-compliance}

The scribe role rests on two bodies of U.S. regulation that oblige financial
institutions to produce specific prose. We cite the primary sources because
the paper's most durable claim---that narration is a task category no tree can
enter---is a legal fact before it is a modeling fact.

\subsection{SAR Narratives (Bank Secrecy Act / FinCEN)}
\label{sec:p8-sar}

Banks must report suspicious transactions to the Financial Crimes Enforcement
Network (FinCEN): under 31~C.F.R.~\S~1020.320, a bank must file a Suspicious
Activity Report for transactions of at least \$5{,}000 conducted at or through
the bank that it ``knows, suspects, or has reason to suspect'' involve illicit
funds, are designed to evade Bank Secrecy Act requirements, or have no
apparent lawful purpose, generally within 30 days of detection
\citep{cfr1020_320}. The report's decisive free-text component is the
narrative. FinCEN's guidance on preparing SAR narratives instructs filers to
identify the essential elements of the suspicious activity---who, what, when,
where, and why, along with the method of operation---in a complete,
self-contained account, and emphasizes that the narrative is where the filing
institution explains why the activity is suspicious
\citep{fincen2003sarnarrative}. This is precisely a conditional
text-generation task: given a case file, produce a complete, accurate,
element-covering narrative. Analyst backlogs make it a bottleneck; its
element-coverage structure makes it checkable; and emerging agentic systems
target exactly this drafting task \citep{naik2025coinvestigator}. Our
architecture keeps the human filer's judgment and signature---the LLM drafts,
the analyst verifies against the evidence links, and the institution files.

\subsection{Adverse-Action Notices (ECOA / Regulation B)}
\label{sec:p8-adverse}

When a creditor takes adverse action on a credit application, the Equal Credit
Opportunity Act as implemented by Regulation~B requires notification including
a statement of specific reasons (or disclosure of the right to obtain one):
under 12~C.F.R.~\S~1002.9, the statement of reasons ``must be specific and
indicate the principal reason(s) for the adverse action,'' and generic appeals
to internal standards or a failing score, without more, do not suffice
\citep{cfr1002_9}. The Consumer Financial Protection Bureau's Circular
2022-03 addresses the machine-learning case directly: the adverse-action
requirements apply regardless of the technology used, creditors may not rely
on ``black-box'' complexity as a defense, and a creditor that cannot state the
specific principal reasons behind its own model's decision is not excused but
non-compliant \citep{cfpb2022circular}.

Two design consequences follow for fraud-adjacent credit decisions. First, the
\emph{scorer} must remain a model whose principal reasons are extractable---a
gradient-boosted tree with feature attributions satisfies this far more
directly than an end-to-end LLM judgment whose verbalized rationale need not
correspond to its computation. Keeping XGBoost in the scorer seat is thus a
compliance posture, not only a performance one. Second, the \emph{scribe} must
be constrained to narrate the scorer's actual attributions: the LLM formats
and phrases the principal reasons; it must not invent them. The draft is
generated from the attribution vector, carries the attribution as provenance,
and is reviewed by a human before issuance. An LLM that freely composed
plausible-sounding reasons would be a Regulation~B violation generator; an LLM
that renders the tree's true top attributions into the required specific
statement is a compliance tool.

\paragraph{Scope note.} We cite U.S. federal requirements because they are the
sharpest published text-generation obligations in fraud and credit workflows;
analogous obligations exist in other jurisdictions but are not analyzed here.
Nothing in this section is legal advice; it establishes only that the
narration category exists, is mandatory, and is textual---hence enterable by
an LLM and not by a tree.
